% ==============================================================
% Combinatorics (Olympiad)
% ==============================================================

\section{Combinatorics (Olympiad)}

\subsection{Principle of inclusion-exclusion}

\begin{theorem}[Inclusion-Exclusion Principle]
For finite sets $A_1, A_2, \ldots, A_n$:
\[
\left|A_1 \cup A_2 \cup \cdots \cup A_n\right|
= \sum_i |A_i| - \sum_{i<j}|A_i \cap A_j| + \sum_{i<j<k}|A_i \cap A_j \cap A_k| - \cdots
\]
For two sets this simplifies to $|A \cup B| = |A| + |B| - |A \cap B|$.
\end{theorem}

\subsubsection*{Summing Multiples with Inclusion-Exclusion}

A standard application is summing all multiples of $a$ or $b$ up to some limit.

\begin{example}[Sum of multiples of 3 or 5 below 1000]
Find the sum of all natural numbers below 1000 that are multiples of 3 or 5.

\textbf{Multiples of 3 below 1000}: these are $3, 6, 9, \ldots, 999$. There are $\lfloor 999/3 \rfloor = 333$ terms, so their sum is
\[
3(1+2+\cdots+333) = 3 \cdot \frac{333 \cdot 334}{2} = 166833.
\]

\textbf{Multiples of 5 below 1000}: $5, 10, \ldots, 995$. There are $\lfloor 999/5 \rfloor = 199$ terms, so their sum is
\[
5(1+2+\cdots+199) = 5 \cdot \frac{199 \cdot 200}{2} = 99500.
\]

\textbf{Multiples of both} (i.e.\ multiples of $\text{lcm}(3,5)=15$): $15, 30, \ldots, 990$. There are $\lfloor 999/15 \rfloor = 66$ terms, so their sum is
\[
15(1+2+\cdots+66) = 15 \cdot \frac{66 \cdot 67}{2} = 33165.
\]

By inclusion-exclusion:
\[
166833 + 99500 - 33165 = \mathbf{233168}.
\]
\end{example}

\begin{remark}[General template]
To sum all multiples of $d$ below $N$: let $k = \lfloor (N-1)/d \rfloor$. The sum is $d \cdot k(k+1)/2$. For multiples of $a$ or $b$ below $N$, subtract the multiples of $\text{lcm}(a,b)$.
\end{remark}

\subsection{Generating functions}

% TODO

\subsection{Bijections and double counting}

% TODO

\subsection{Combinatorial game theory}

\subsubsection*{Two-Pile Threshold Games}

Consider the following two-player game: two piles of stones of sizes $a$ and $b$ ($a \le b$). A move consists of removing at least one stone from each pile, with the \emph{total} number removed equal to $\min(a,b)$. The player who cannot move loses.

\textbf{Key reduction.} From a position $(a,b)$ with $a \le b$, exactly $a$ stones are removed in total (since $a = \min(a,b)$). The total remaining is $(a+b)-a = b$. So every move from $(a,b)$ lands on a position whose pile sum is exactly $b$. Equivalently, the reachable positions from $(a,b)$ are exactly:
\[
\{(c,\, b-c) : 1 \le c \le \min(a, \lfloor b/2 \rfloor)\}.
\]
For fixed $b$, increasing $a$ only \emph{adds} reachable moves — it never removes them. This gives a monotonicity principle:

\begin{theorem}[Column threshold]
For fixed $b$, there exists a threshold $r(b)$ such that $(a,b)$ is a losing position if and only if $a < r(b)$.
\end{theorem}

\subsubsection*{The Lowbit Threshold}

\begin{theorem}
The threshold for column $b$ is
\[
r(b) = \operatorname{lowbit}(b+1) = 2^{v_2(b+1)},
\]
i.e.\ the largest power of $2$ dividing $b+1$. The position $(a,b)$ (with $a \le b$) is \textbf{losing} if and only if
\[
a < \operatorname{lowbit}(b+1).
\]
\end{theorem}

\begin{proof}[Proof sketch]
Let $r = \operatorname{lowbit}(b+1)$. A child $(c, b-c)$ is losing iff $c < \operatorname{lowbit}(b-c+1)$.

\emph{Case $1 \le c < r$}: Then $\operatorname{lowbit}((b+1)-c) = \operatorname{lowbit}(c) \le c$, so the child is \emph{not} losing (it is a winning position).

\emph{Case $c = r$}: Since $r = \operatorname{lowbit}(b+1)$, subtracting $r$ from $b+1$ eliminates the lowest set bit, making $\operatorname{lowbit}(b+1-r) > r$. So $r < \operatorname{lowbit}(b-r+1)$, meaning $(r, b-r)$ \emph{is} losing.

Therefore: the first losing child is at $c = r$. Any position with $a < r$ cannot reach it, so $(a,b)$ is losing. Any position with $a \ge r$ can move to $(r, b-r)$ (losing for the opponent), so it is winning.
\end{proof}

\begin{example}[Small cases]
\begin{center}
\begin{tabular}{ccll}
$b$ & $b+1$ & $\operatorname{lowbit}(b+1)$ & Losing positions $(a,b)$, $a \le b$ \\
\hline
3 & 4 & 4 & all of $(1,3),(2,3),(3,3)$ \\
5 & 6 & 2 & only $(1,5)$ \\
7 & 8 & 8 & all of $(1,7),\ldots,(7,7)$ \\
11 & 12 & 4 & $(1,11),(2,11),(3,11)$ \\
\end{tabular}
\end{center}
The ``all-losing columns'' occur when $\operatorname{lowbit}(b+1) > b$, i.e.\ when $b+1$ is itself a power of 2, so $b = 2^k - 1$: columns $1, 3, 7, 15, 31, \ldots$
\end{example}

\subsubsection*{Counting Losing Positions}

\begin{theorem}[Number of ordered losing pairs $L(n)$]
For ordered pairs $(a,b)$ with $1 \le a,b \le n$:
\[
L(n) = 2\bigl(T(n+1) - 1\bigr) - 2n - \lfloor \log_2(n+1) \rfloor,
\]
where $T(N) = \sum_{m=1}^{N} \operatorname{lowbit}(m)$.
\end{theorem}

\begin{proof}[Derivation]
For each column $b$, the number of losing values $a \in \{1,\ldots,b\}$ is $\operatorname{lowbit}(b+1) - 1$. For ordered pairs, off-diagonal positions $(a,b)$ and $(b,a)$ are counted separately; the diagonal $(b,b)$ is losing exactly when $b < \operatorname{lowbit}(b+1)$, i.e.\ when $b = 2^k - 1$. Summing and adjusting for symmetry gives the formula above.
\end{proof}

\begin{remark}[Fast computation]
The sum $T(N)$ satisfies the recurrence $T(2k) = k + 2T(k)$ and $T(2k+1) = T(2k)+1$ (see the 2-adic valuation subsection in the number theory chapter), enabling $O(\log N)$ evaluation. This makes $L(n)$ computable in logarithmic time for huge $n$.
\end{remark}

\begin{example}
$L(7) = 21$, $L(49) = 221$, $L(7^{17}) = 10784223938983273$.
\end{example}
